\documentclass[10pt,a4paper]{amsart}
\usepackage[margin=19mm]{geometry}
\usepackage{amsmath,amssymb,mathtools}
\usepackage[scr=rsfs,cal=euler]{mathalfa}
\usepackage{xfrac}
\usepackage[unicode,hidelinks,bookmarks=false]{hyperref}
\newcommand{\ie}{i.e.\ }
\newcommand{\cf}{cf.\ }
\newcommand{\resp}{respectively\ }
\newcommand{\Subsection}{\subsection}
\providecommand{\nobreakdash}{\nobreak\hbox{-}}
\setlength{\emergencystretch}{2em}
\multlinegap0pt
\usepackage{amssymb}
\usepackage[shortlabels]{enumitem}
\usepackage{float}
\newtheorem{lemma}{Lemma}
\newcommand{\CZ}{\mu_{CZ}^{\tau}}
\newcommand{\Q}{\mathbb{Q}}
\newcommand{\R}{\mathbb{R}}
\newcommand{\ind}{\textnormal{ind}}
\title{Cylindrical contact homology for weakly convex contact forms in dimension three}
\author{Ana Kelly de Oliveira, Pedro A. S. Salom\~ao}
\date{Source: arXiv:2603.29352v1 (CC BY 4.0)}
\begin{document}
\maketitle

\noindent\emph{Source and licence.} Cylindrical contact homology for weakly convex contact forms in dimension three. Version arXiv:2603.29352v1 (CC BY 4.0), distributed by arXiv under the Creative Commons Attribution 4.0 International licence, http://creativecommons.org/licenses/by/4.0/, as recorded in the arXiv licence metadata for this exact version and shown on the abstract page https://arxiv.org/abs/2603.29352v1. Full licence notice: \texttt{licenses/arxiv-2603.29352-CC-BY-4.0.txt}.\par\smallskip

\noindent\textit{Editorial scope.} Counted unit: the article's lemma odepeq (source0.tex lines 1332--1340). The article's complete original proof of it is delivered with it (source0.tex lines 1341--1376, one \texttt{proof} environment of 36 lines). No mathematics is written, repaired or re-derived: the statement and the proof are the article's own lines, in the article's own notation, and no retained line is substituted or re-flowed.  No further source-local context is needed: every cross-reference of every retained line resolves inside the two delivered spans.  Nothing was omitted from any retained span, so the package carries no omission record.  No retained line contains a citation, so the wrapper carries no bibliography and no bibliography entry.  No retained line contains a figure environment, an image-inclusion command or drawing code, so no image is delivered and the package declares no asset dependency.  Editorial formatting only, in addition: an AMS \texttt{amsart} preamble that restates the article's own declarations and macros used by the retained lines, namely the environments \texttt{lemma} on separate counters, together with the standard packages those lines need.  The retained environments print in this document as Lemma 1 in order of appearance, whereas the article prints the same items with the numbering of its own full document.  The complete native source of the article as distributed in the arXiv e-print for this exact version is bundled unchanged as \texttt{sources/197510/source0.tex}.
\begin{lemma}\label{odepeq}
Assume that $\alpha_1^d, \alpha_2^d$ are good orbits satisfying $\CZ(\alpha_1^d)-\CZ(\alpha_2^d)=2$. Then the following assertions hold:
\begin{itemize}
    \item[(i)] If $\alpha_1$ and $\alpha_2$ are both hyperbolic, then $d=1$.
    \item[(ii)] If $\alpha_1$ is elliptic and $\alpha_2$ is hyperbolic, then $\alpha_2$ is negative hyperbolic and $d=1$.
    \item[(iii)] If $\alpha_1$ is hyperbolic and $\alpha_2$ is elliptic, then $\alpha_1$ is negative hyperbolic and $d=1$.
\end{itemize}
In particular, if $d>1$, then $\alpha_1$ and $\alpha_2$ are elliptic.
\end{lemma}

\begin{proof}
We have
\begin{eqnarray*}
1&=&\ind(\mathfrak{v}_1)=-(2-3)+\CZ(\alpha_1)-\CZ(\alpha_2)-2\\
&=& \CZ(\alpha_1)-\CZ(\alpha_2)-1,
\end{eqnarray*}
which implies
$\CZ(\alpha_1)-\CZ(\alpha_2)=2.$
In particular, $\CZ(\alpha_1)$ and $\CZ(\alpha_2)$ have the same parity. Therefore,
\begin{itemize}
    \item If $\alpha_1$ and $\alpha_2$ are hyperbolic, then our hypotheses imply
     $$2=\CZ(\alpha_1^d)-\CZ(\alpha_2^d)=d(\CZ(\alpha_1)-\CZ(\alpha_2))=2d \,\,\Rightarrow \,\, d =1.$$
     \item If $\alpha_1$ is elliptic and $\alpha_2$ is hyperbolic, then we obtain
     \begin{eqnarray*}
     2&=& \CZ(\alpha_1^d)-\CZ(\alpha_2^d)=2\lfloor d\theta_{\alpha_1}\rfloor+1 -d\cdot\CZ(\alpha_2)\\
     &\geq & 2d\lfloor \theta_{\alpha_1}\rfloor+1 -d \cdot \CZ(\alpha_2)=d(2\lfloor \theta_{\alpha_1}\rfloor+1)-d+1 -d \cdot\CZ(\alpha_2)\\
     &=&d(\CZ(\alpha_1)-\CZ(\alpha_2))-d+1\\
     &=& d+1.
     \end{eqnarray*}
     Hence  $d=1$.
     \item Assume now that $\alpha_1$ is hyperbolic and $\alpha_2$ is elliptic. Let $\{x\}:=x-\lfloor x \rfloor, \forall x\in \R$. Then $\{\theta_{\alpha_2}\} \in [0,1] \backslash \Q$, which implies $0<d\{\theta_{\alpha_2}\}<d$. We estimate
    \begin{eqnarray*}
    \lceil d\theta_{\alpha_2} \rceil &=& \big\lceil d \lfloor \theta_{\alpha_2} \rfloor + d\{\theta_{\alpha_2}\} \big\rceil = d \lfloor \theta_{\alpha_2} \rfloor + \lceil d\{\theta_{\alpha_2}\} \rceil \\
    &\leq & d \lfloor \theta_{\alpha_2} \rfloor +d = d(\lfloor \theta_{\alpha_2} \rfloor + 1) \\
    &=& d\lceil \theta_{\alpha_2}\rceil,
    \end{eqnarray*}
    and
    \begin{eqnarray*}
    2&=& \CZ(\alpha_1^d)-\CZ(\alpha_2^d)= d \cdot \CZ(\alpha_1)-(2\lceil d\theta_{\alpha_2}\rceil-1)\\
    &\geq& d \cdot \CZ(\alpha_1)-2d \lceil \theta_{\alpha_2} \rceil + 1 = d \cdot \CZ(\alpha_1)-d(2 \lceil \theta_{\alpha_2} \rceil -1)-d+1 \\
    &=& d(\CZ(\alpha_1)-\CZ(\alpha_2))-d+1. \\
    &=& d+1
    \end{eqnarray*}
    We conclude as before that $d=1$.
\end{itemize}
\end{proof}

\end{document}
